 %!TEX root = main.tex

\subsection*{Network Overview}

GlobalFlow is an international third-party logistics operator managing the
end-to-end distribution of goods across a directed network
$G=(N,A)$ with a four-tier hierarchical structure.
The network connects $|S|=9$ origin suppliers to $|C|=30$ final customers
through $|H|=5$ international hubs and $|W|=15$ regional warehouses,
for a total of $|N|=59$ nodes and $|A|=317$ arcs.

\vspace{4pt}
{\small
\begin{tabularx}{\columnwidth}{@{}l X r@{}}
\toprule
\textbf{Tier} & \textbf{Key locations} & $|\cdot|$ \\
\midrule
Suppliers ($S$)   & 3 per product (see below)                  & 9  \\
Intl.\ Hubs ($H$) & FRA, DXB, SIN, ORD, GRU                   & 5  \\
Warehouses ($W$)  & Rotterdam, Madrid, Mumbai \ldots            & 15 \\
Customers ($C$)   & Europe~8, USA/Canada~8, Brazil~5, Japan~5, India~4 & 30 \\
\bottomrule
\end{tabularx}
}

\subsection*{Products}
 
Three product families share the network infrastructure but differ in their
cost structure, supplier base, and disruption exposure:
 
\begin{description}[
  labelwidth=1.2em, leftmargin=1.8em,
  itemsep=2pt, parsep=0pt, font=\bfseries\normalsize]
  \item[A] \textit{Fertilisers} : S1–S3 (Minsk, Casablanca, Port of Spain);
    sea and road transport; demand concentrated in agricultural regions.
  \item[B] \textit{Semiconductors} : S4–S6 (Taipei, Seoul, Yokohama);
    long-haul sea freight with road last-mile delivery; heavily exposed
    to East Asian hub disruptions.
  \item[C] \textit{Battery components} : S7–S9 (Antofagasta, Lubumbashi, Chengdu);
    sea and road transport; sensitive to both tariff and energy shocks.
\end{description}


\subsection*{Arcs and Cost Structure}

Of the 317 arcs, 270 are \textit{always-active} ($A_{\text{always}}$, zero
fixed cost) and 47 are \textit{optional} ($A_{\text{fixed}}$, positive
activation cost). Each arc carries a shared capacity $u_a$ across all products.
Variable costs reflect distance, modal choice (air $\gg$ sea $>$ road), and
product characteristics. Inter-zone tariff rates are applied multiplicatively
on top of the baseline variable cost according to the bilateral zone matrix in
\texttt{globalflow\_instance.xlsx}.

\paragraph*{Problem classification.}
The GlobalFlow design problem belongs to the class of
Multi-Commodity Capacitated Fixed-Charge Network Design Problems (MC-CFCNDP).
It is NP-hard in general: the binary facility-location
and arc-activation decisions interact with continuous multi-commodity
flows through shared-capacity constraints, making the LP relaxation
generally fractional.
The practical tractability of this instance ($|N|=59$, $|A|=317$,
$|P|=3$) is confirmed empirically in Section~\ref{sec:phase1}:
FICO Xpress solves the MIP to global optimality in under $0.1$\,s,
reflecting a well-structured instance with tight LP–IP gaps.

\begin{tcolorbox}[questionbox, title={Central Question}]
\small\itshape
Is a network optimal under normal operating conditions also robust to the
geopolitical and economic shocks of today's environment?
And if not, at what cost can resilience be achieved?
\end{tcolorbox}
